\documentclass[11pt]{article}
\usepackage[margin=1in]{geometry}
\usepackage{amsmath,amssymb,amsthm,hyperref}
\newtheorem{theorem}{Theorem}
\newcommand{\ess}{\mathrm{ess}}
\newcommand{\sigp}{\sigma_{\mathrm p}}
\title{The weighted Dirichlet operator has infinitely many eigenvalues\\
\large MathDB 366735}
\author{Principia Math in Progress}
\date{12 September 2026}
\begin{document}
\maketitle

\paragraph{Task 1.}
Resolve Conjecture 6.1 of Monsalve-L\'opez and Seco \cite{MS}.
For every $\alpha>0$ and nonzero $a\in\mathbb C$, it asserts that the operator
\[
 V_a=M_{a-z}^*M_{a-z}
\]
has empty point spectrum on the space with norm
\[
 \left\|\sum_{n\geq0}f_nz^n\right\|_\omega^2
 =\sum_{n\geq0}|f_n|^2\omega_n,\qquad
 \omega_n=\binom{n+\alpha}{n}.
\]
The conjecture is false. A one-vector counterexample suffices, and a
localized variational argument proves the opposite behavior throughout
the stated parameter range.

\begin{theorem}\label{thm:infinite}
For every $\alpha>0$ and $a\ne0$, the operator $V_a$ has infinitely many
isolated eigenvalues of finite multiplicity strictly above $(1+|a|)^2$,
accumulating at $(1+|a|)^2$. Its essential spectrum is
\begin{equation}\label{eq:ess}
 \sigma_{\ess}(V_a)=[(1-|a|)^2,(1+|a|)^2].
\end{equation}
In particular, for $\alpha=1$ and $a=1/4$ there is an eigenvalue
$\lambda\geq33/16>25/16$.
\end{theorem}

\paragraph{Weighted-shift representation.}
The monomials $e_n=z^n/\sqrt{\omega_n}$ form an orthonormal basis.
Writing $S=M_z$, we have
\begin{equation}\label{eq:shift}
 Se_n=q_ne_{n+1},\qquad
 q_n^2=\frac{\omega_{n+1}}{\omega_n}
      =1+\frac{\alpha}{n+1}.
\end{equation}
Hence $S$ is bounded, $q_n>1$, and $q_n\to1$. Moreover,
\begin{equation}\label{eq:operator}
 V_a=|a|^2I-\overline aS-aS^*+S^*S
\end{equation}
is bounded, positive, and self-adjoint.
If $Ue_n=e_{n+1}$, then $S-U$ is compact, since truncating its weight
sequence gives finite-rank operators converging in operator norm.
Likewise $S^*S-I$ is a compact diagonal operator. Thus $V_a$ is a compact
perturbation of
\[
 W_a=(|a|^2+1)I-\overline aU-aU^*.
\]
The free half-line Jacobi operator $U+U^*$ has spectrum and essential
spectrum $[-2,2]$. A diagonal unitary removes the argument of $a$.
Weyl invariance of the essential spectrum proves \eqref{eq:ess},
also given by Theorem 1.2 of \cite{MS}.
Every spectral point outside this interval is an isolated eigenvalue
of finite multiplicity.

\paragraph{An exact counterexample.}
Since $e_0=1$ has norm one and $1,z$ are orthogonal,
\begin{equation}\label{eq:onevector}
 \langle V_a e_0,e_0\rangle
 =\|a-z\|_\omega^2=|a|^2+1+\alpha.
\end{equation}
Its excess over the upper essential edge is $\alpha-2|a|$.
Whenever $\alpha>2|a|$, the supremum of the spectrum therefore lies
strictly above that edge and is an eigenvalue.
For $\alpha=1$, $a=1/4$, the expectation is $33/16$ and the edge is
$25/16$, with exact gap $1/2$. This already refutes the full universal
conjecture.

\paragraph{Infinitely many eigenvalues for all parameters.}
Put $\rho=|a|>0$ and $E=(1+\rho)^2$.
In a suitable diagonal unitary basis, the matrix of $V_a$ has diagonal
$1+\rho^2+\alpha/(n+1)$ and positive off-diagonal entries $\rho q_n$.
For a real nonnegative finitely supported vector $u=(u_n)$ its quadratic
form, denoted $Q(u)$, consequently satisfies
\begin{align}
 Q(u)-E\|u\|_2^2
 &\geq \sum_{n\geq0}\frac{\alpha u_n^2}{n+1}
       -2\rho\sum_{n\geq0}u_n^2
       +2\rho\sum_{n\geq0}u_nu_{n+1}\notag\\
 &=\sum_{n\geq0}\frac{\alpha u_n^2}{n+1}
       -\rho\sum_{n=-1}^{\infty}(u_{n+1}-u_n)^2,
       \qquad u_{-1}=0.                         \label{eq:form}
\end{align}
Here $q_n\geq1$ and nonnegativity of the coordinates justify the first
inequality; zero coordinates are used outside the support.

For an integer $L\geq1$, place a tent vector on $[L,3L-2]$:
\[
 u^{(L)}_{L+j-1}=\min(j,2L-j)\quad(1\leq j\leq2L-1),
 \qquad u^{(L)}_n=0\text{ elsewhere}.
\]
Direct summation gives
\begin{equation}\label{eq:tent}
 N_L:=\|u^{(L)}\|_2^2
 =L^2+2\sum_{j=1}^{L-1}j^2=\frac{2L^3+L}{3},
 \qquad
 \sum_{n=-1}^{\infty}(u^{(L)}_{n+1}-u^{(L)}_n)^2=2L.
\end{equation}
On its support, $n+1\leq3L$. Therefore \eqref{eq:form} yields
\begin{equation}\label{eq:positive}
 \frac{Q(u^{(L)})}{N_L}-E
 \geq\frac{\alpha}{3L}-\frac{6\rho}{2L^2+1}
 >\frac{\alpha}{3L}-\frac{3\rho}{L^2}>0
 \quad\text{if }L>\frac{9\rho}{\alpha}.
\end{equation}

Choose such an $L_0$ and put $L_j=4^jL_0$.
The supports of the corresponding tent vectors are separated by at
least one zero coordinate. The vectors are orthogonal, and the Jacobi
matrix has no quadratic-form cross terms between them.
Thus, for every $k$, the span of the first $k$ vectors is a $k$-dimensional
subspace on which $Q(u)>E\|u\|_2^2$ for all nonzero $u$.

To see the spectral implication directly, if the spectral projection of
$V_a$ onto $(E,\infty)$ had rank less than $k$, that subspace would contain
a nonzero vector orthogonal to its range. The spectral theorem would
give $Q(u)\leq E\|u\|_2^2$, a contradiction.
The projection therefore has infinite rank. Since $(E,\infty)$ is
outside the essential spectrum, it consists of isolated eigenvalues
of finite multiplicity. There must be infinitely many of them.
Boundedness and the absence of accumulation outside the essential
spectrum force their only accumulation point to be $E$.
This proves Theorem \ref{thm:infinite}.
\hfill$\square$

\paragraph{Scope and verification.}
The norm is exactly the one in the published conjecture. The paper's
Theorem 1.5 treats a different, equivalent norm for $0<\alpha<1$;
it does not assert this result for the displayed binomial weights.
The general argument above is independent of finite matrix truncations.
Exact checks verify the weight and coefficient normalization, the
one-vector certificate, and the tent identities. Supplementary interval
checks evaluate the actual Jacobi forms of selected tents.

\paragraph{Task 2.}
Determine the leading asymptotic decay of the eigenvalues above
$(1+|a|)^2$, including its dependence on $\alpha$ and $|a|$, and obtain
an effective error bound.

\paragraph{Task 3.}
Classify the eigenvalues below $(1-|a|)^2$ and determine how their
existence and accumulation change as $|a|$ crosses one.

\paragraph{Task 4.}
Find sharp tail conditions on positive weights $\omega_n$ that ensure
infinitely many eigenvalues above the free upper edge. Analyze the
critical case $\omega_{n+1}/\omega_n-1$ of order $n^{-2}$.

\paragraph{Task 5.}
Determine whether embedded or edge eigenvalues occur for the binomial
weights, and describe the absolutely continuous and singular continuous
parts of the spectral measure.

\begin{thebibliography}{9}
\bibitem{MS}
M. Monsalve-L\'opez and D. Seco,
\emph{Towards spectral descriptions of cyclic functions},
Annals of Functional Analysis \textbf{16} (2025), article 52.
Theorem 1.2, Section 2, and Conjecture 6.1.
\href{https://doi.org/10.1007/s43034-025-00446-0}
{doi:10.1007/s43034-025-00446-0}.
\end{thebibliography}
\end{document}
